\section*{الفصل \ref{ch:b3:remote-sensing-gis} --- الاستشعار عن بعد و GIS}

\begin{solution}{exo:b3:remote-sensing-gis:1}
الحقل: $(0.42 - 0.06)/(0.42 + 0.06) = 0.75$، غطاء نباتي أخضر كثيف، لمحصول
في أوج اكتساء أوراقه. الحقل المحروث: $(0.28 - 0.21)/(0.28 + 0.21) = 0.14$، أرض شبه عارية،
تربة تزداد سطوعًا قليلًا نحو الأشعة تحت الحمراء.
\end{solution}

\begin{solution}{exo:b3:remote-sensing-gis:2}
(أ)~$185000/30 \approx 6170$ بكسلًا؛ و$185000/15 \approx 12300$
في الحزمة البانكروماتية. (ب)~$2^{12} = 4096$ مستوى؛ $2^8 = 256$. (ج)~16 يومًا
منفردًا؛ ومع لاندسات~9 في المدار نفسه بفارق 8 أيام، 8 أيام.
\end{solution}

\begin{solution}{exo:b3:remote-sensing-gis:3}
(أ)~الهدب الواحد هو $\lambda/2 = \qty{27.7}{mm}$؛ وستة أهداب تساوي
\qty{166}{mm} من التغيّر على \omterm{def:b3:remote-sensing-gis:insar}{خط النظر}، وهو هنا مبتعد عن القمر الصناعي، لأنّ
الحوض يهبط. (ب)~$56.6/2 = \qty{28.3}{mm}$.
\end{solution}

\begin{solution}{exo:b3:remote-sensing-gis:4}
(أ)~$t = 2d/c = 2\times\num{5e5}/\num{3.00e8} = \qty{3.34}{ms}$.
(ب)~$c\,\delta t/2 = \num{3.00e8}\times\num{8e-10}/2 = \qty{0.12}{m}$؛ إنّ سطحًا
يُحدَّد موقع سطحٍ تُقاس منه فوتونات كثيرة بدقّة أفضل بكثير، وتتناقص الدقّة
كمقلوب الجذر التربيعي لعددها.
\end{solution}

\begin{solution}{exo:b3:remote-sensing-gis:5}
(أ)~بوحدات النظام الدولي $L = \qty{9.0e6}{W.m^{-2}.sr^{-1}.m^{-1}}$. ومع $c_1 =
2hc^2 = \qty{1.191e-16}{W.m^2.sr^{-1}}$، $c_2 = \qty{0.014388}{m.K}$
و$\lambda^5 = (\num{1.09e-5})^5 = \num{1.539e-25}$، $c_1/(\lambda^5L) = 86.0$،
و$T_b = c_2/(\lambda\ln 87.0) = \qty{295.6}{K}$. (ب)~درجة الحرارة الحركية
تحقّق $0.98\,B_\lambda(T) = L$، أي $B_\lambda(T) = L/0.98$: تعطي الصيغة نفسها
$T = \qty{296.9}{K}$، أي بزيادة \qty{1.3}{K}.
\end{solution}

\begin{solution}{exo:b3:remote-sensing-gis:6}
(أ)~$R = 693/\cos\ang{35} = \qty{846}{km}$. (ب)~الهوائي الحقيقي: $R\lambda/L =
\num{8.46e5}\times0.0555/12.3 = \qty{3.8}{km}$؛ فتحة اصطناعية: $L/2 =
\qty{6.2}{m}$، أي نحو 600 مرة أدقّ. (ج)~الفتحة الاصطناعية هي طول المسار
الذي تبقى خلاله نقطة في الحزمة، $R\lambda/L =
\qty{3.8}{km}$.
\end{solution}

\begin{solution}{exo:b3:remote-sensing-gis:7}
$\partial z/\partial x = (418 - 405)/60 = 0.217$ (الشرق ناقص الغرب)،
$\partial z/\partial y = (418 - 401)/60 = 0.283$ (الشمال ناقص الجنوب)؛
$|\nabla z| = 0.357$، ميلًا قدره $\arctan 0.357 = \ang{19.6}$. ترتفع الأرض
نحو الشمال الشرقي، فتواجه الجنوب الغربي: $A =
\operatorname{atan2}(-0.217, -0.283) = \ang{217}$. التظليل الجبلي: $\cos\ang{45}
\cos\ang{19.6} + \sin\ang{45}\sin\ang{19.6}\cos(\ang{315} - \ang{217}) =
0.666 - 0.033 = 0.63$، أغمق قليلًا من الأرض المستوية (0.71): يبتعد الميل
ينحرف قليلًا بعيدًا عن مصدر ضوء في الشمال الغربي.
\end{solution}

\begin{solution}{exo:b3:remote-sensing-gis:8}
مجاميع الصفوف 52، 48، 38، 22؛ ومجاميع الأعمدة 54، 46، 38، 22؛ القطر 133 من
160: الدقّة الكلّية \qty{83.1}{\percent}. دقّة المستخدم: المياه
$48/52 = 0.92$، الغابات $37/48 = 0.77$، المحاصيل $29/38 = 0.76$، المدينة
$19/22 = 0.86$. دقّة المنتج: المياه $48/54 = 0.89$، الغابات $37/46 = 0.80$،
المحاصيل $29/38 = 0.76$، المدينة $19/22 = 0.86$. $p_e = (52\cdot54 + 48\cdot46 +
38\cdot38 + 22\cdot22)/160^2 = 0.271$، $\kappa = (0.831 - 0.271)/(1 -
0.271) = 0.77$. تُغفَل المحاصيل أكثر من غيرها: ذهب 9 من أصل 38 بكسلًا
محصوليًا حقيقيًا إلى أماكن أخرى، معظمها إلى الغابات.
\end{solution}

\begin{solution}{exo:b3:remote-sensing-gis:9}
(أ)~$L = 0.04\,B_\lambda(1443) + 0.96\,B_\lambda(290)$ عند
\qty{10.9}{\micro\metre}، \qty{28.6}{W.m^{-2}.sr^{-1}.\micro\metre^{-1}}،
تبلغ درجة حرارة سطوعه \qty{396}{K}، في مقابل متوسط درجة الحرارة
$0.04\times1443 + 0.96\times290 = \qty{336}{K}$. (ب)~عند
\qty{2.2}{\micro\metre}، $B_\lambda(1443)/B_\lambda(290) = \num{6.8e7}$:
تهيمن الحمم على البكسل تمامًا، وتبلغ درجة حرارة سطوعها
\qty{845}{K}. تكشف حزمة الأشعة تحت الحمراء القصيرة البقع الساخنة الصغيرة
على نحو أفضل بكثير؛ أمّا الحزمة الحرارية، الحسّاسة للبكسل كلّه، فهي التي
تقيس الأرض المحيطة بها.
\end{solution}

\begin{solution}{exo:b3:remote-sensing-gis:10}
مع وجود القمر الصناعي إلى الشرق، $d_1 = u_E\sin\theta + u_z\cos\theta$؛ وإلى
الغرب، $d_2 = -u_E\sin\theta + u_z\cos\theta$. بالجمع والطرح،
$u_z = (d_1 + d_2)/(2\cos\theta) = 102/1.638 = \qty{62.3}{mm}$ و$u_E =
(d_1 - d_2)/(2\sin\theta) = 26/1.147 = \qty{22.7}{mm}$. يعطي مسار المسح الواحد معادلةً واحدة
لمجهولين: أيّ مزيج من الرفع والحركة نحو الشرق بالإسقاط نفسه على \omterm{def:b3:remote-sensing-gis:insar}{خط النظر}
يعطي الأهداب نفسها.
\end{solution}

\begin{solution}{exo:b3:remote-sensing-gis:11}
(أ)~حسب \cref{thm:b3:remote-sensing-gis:fringe}، $\Delta\phi =
4\pi\Delta r/\lambda < \pi$ يتطلّب $\Delta r < \lambda/4 =
\qty{13.9}{mm}$. (ب)~$u_z = k z_0(r^2 + z_0^2)^{-3/2}$ مع $k = (1 -
\nu)\Delta V/\pi$; $\mathrm du_z/\mathrm dr = -3kz_0r(r^2 + z_0^2)^{-5/2}$،
التي تنعدم مشتقّتها عندما $(r^2 + z_0^2) = 5r^2$، أي عند $r = z_0/2$. وعندها
\[
\left|\frac{\mathrm du_z}{\mathrm dr}\right| = \frac{3k}{2z_0^3}\Big(\frac54\Big)^{-5/2}
= \frac{0.859\,k}{z_0^3} = \frac{0.859\times\num{2.39e6}}{\num{6.4e10}} = \num{3.2e-5},
\]
\qty{32}{mm} لكلّ كيلومتر. (ج)~$13.9/32 = \qty{0.43}{km}$: بكسلات أصغر
من نحو \qty{430}{m}، وهو ما يتحقّق بسهولة مع سنتينل-1.
\end{solution}

\begin{solution}{exo:b3:remote-sensing-gis:12}
(أ)~من أجل $N, R \ge 0$ ليسا معًا صفرًا، $|N - R| \le N + R$، ومنه $|\mathrm{NDVI}|
\le 1$؛ أمّا استبدال $N$ و$R$ بالمقدارين $aN$ و$aR$ فيضرب البسط والمقام
في المقام بالمقدار $a$. (ب)~بدون التصحيح: $(0.502 - 0.051)/(0.502 + 0.051) =
0.816$؛ ومع التصحيح: $(0.512 - 0.081)/(0.512 + 0.081) = 0.727$. (ج)~يخفّض إشعاع المسار
مؤشّر الغطاء النباتي الأخضر بنحو العُشر هنا،
ويتغيّر هذا المقدار مع الضباب: ستُظهر خريطة الصور غير المصحّحة
الطقسَ بقدر ما تُظهر الغطاءَ النباتي، وتصبح المقارنات بين التواريخ
عديمةَ المعنى.
\end{solution}

\begin{solution}{pb:b3:remote-sensing-gis:1}
\textbf{1.} الغابة $(0.38 - 0.04)/0.42 = 0.81$؛ تدفّق عام 1890 $0.23/0.37 =
0.62$؛ تدفّق عام 1960 $0.10/0.32 = 0.31$؛ تدفّق عام 2005 $0.02/0.26 = 0.08$.

\textbf{2.} $1890$، ثم 1960، ثم 2005: أقدمُ التدفّقات هو الأكثرُ خضرةً،
وفق ترتيب الأعمار؛ ولم يصل بعدُ إلى الغابة.

\textbf{3.} تعكس الصخور العارية قليلًا أكثر في تحت الأحمر القريب منها في
الأحمر، مثل الرمل البنّي في شكل الفصل؛ وتضيف الأشنات وأولى النباتات
قليلًا: 0.08 هي قيمة الأرض شبه العارية.

\textbf{4.} يُخفِّت المنحدرُ المُولّى عن الضوء الحزمتين بالنسبة نفسها،
فتُلغيه النسبة، بينما يخلط الانعكاس تحت الأحمر القريب وحده
بين الغابة المظلَّلة والغطاء النباتي المتفرّق.

\textbf{5.} حسب \cref{prop:b3:remote-sensing-gis:planck-inversion} مع
$\lambda = \qty{10.9}{\micro\metre}$: الهضبة \qty{291.1}{K}، وبكسل
بكسل الفومارول \qty{302.3}{K}.

\textbf{6.} $0.97\,B_\lambda(T) = \num{8.40e6}$ يعطي $T = \qty{293.0}{K}$.

\textbf{7.} $B_\lambda(\qty{370}{K}) =
\qty{22.48}{W.m^{-2}.sr^{-1}.\micro\metre^{-1}}$; $f = (9.95 - 8.40)/(22.48
- 8.40) = 0.110$، نحو \qty{1100}{m^2} من البكسل ذي المساحة \qty{e4}{m^2}.

\textbf{8.} الإشعاع دالّة محدّبة في درجة الحرارة
(\cref{prop:b3:remote-sensing-gis:mixed}): الجزء الحار يضيف إشعاعًا أكبر
من حصّته في درجة الحرارة، فيبدو البكسل أسخن من متوسّطه،
$0.110\times370 + 0.890\times291.1 = \qty{299.8}{K}$ مقابل
\qty{302.3}{K}.

\textbf{9.} يتوقّف طور كلّ صدى على المسار الدقيق؛ فتغيير كبير في زاوية الرصد
يغيّر مجموع أصداء المبعثرات، ويفقد
مخطّط التداخل تماسكه، ويضيف طورًا من التضاريس.
$48/12 = 4$ دورات لقمر صناعي واحد.

\textbf{10.} $5.4\times27.73 = \qty{149.7}{mm}$ باتجاه القمر الصناعي.

\textbf{11.} $u_z = 149.7/\cos\ang{35} = \qty{182.8}{mm}$.

\textbf{12.} تتحرّك الأرض نحو الخارج كما ترتفع؛ وعلى الجانب الشرقي
تتّجه الحركة الخارجية نحو القمر الصناعي فتُضاف إلى الرفع، وعلى الجانب الغربي
تُطرح منه (\cref{prop:b3:remote-sensing-gis:mogi-los}). ويتناقص التغيّر
أسرع نحو الغرب، حيث تتزاحم الأهداب؛ وينمدّ النمط
نحو الشرق.

\textbf{13.} على ذلك الخط تكون الحركة الأفقية شمال--جنوب،
عمودية على \omterm{def:b3:remote-sensing-gis:insar}{خط النظر}، فيرى الرادار $u_z\cos\theta$ وحده، بشكل
الرفع نفسه. ينصف الرفع عند $r = 0.766\,z_0$، فيكون $z_0 =
3.4/0.766 = \qty{4.44}{km}$.

\textbf{14.} نصف 5.4: 2.7 أهداب.

\textbf{15.} يكون التشوّه متطابقًا في مخطّطَي التداخل بعد
تعديل الفواصل الزمنية، أو يستمرّ في الازدياد من زوج إلى زوج؛ أمّا
النمط الجوي فيتغيّر من تاريخ إلى آخر، وغالبًا ما يتّبع
التضاريس. ويؤكّد التشوّهَ زوجٌ ثانٍ يُظهر
نمط عين الثور نفسه.

\textbf{16.} $\Delta V = \pi u_z(0)z_0^2/(1 - \nu) = \pi\times0.1828\times
4436^2/0.75 = \qty{1.51e7}{m^3}$.

\textbf{17.} $\num{1.51e7}/48 = \qty{3.1e5}{m^3}$ في اليوم، \qty{1.1e8}{m^3}
في السنة إن بقي المعدّل ثابتًا.

\textbf{18.} $u_r \propto r/(r^2 + z_0^2)^{3/2}$ يبلغ أكبر قيمة له حيث تنعدم
مشتقّته، $r^2 + z_0^2 = 3r^2$، $r = z_0/\sqrt2 = \qty{3.14}{km}$؛
وعندها $u_r = \qty{70}{mm}$، أي \qty{38}{\percent} من الرفع المركزي.

\textbf{19.} نصف القطر \qty{1.48}{km}، الحجم $\frac43\pi(1.48)^3 =
\qty{13.6}{km^3}$؛ والتغيّر هو $\num{1.51e7}/\num{1.36e10} =
\qty{0.11}{\percent}$ منه.

\textbf{20.} لا: عُشر الهدب يساوي \qty{2}{\percent} من العدّ،
والعمقُ، المقروءُ من نصف عرضٍ يبلغ \qty{3.4}{km} ربّما إلى
\qty{0.1}{km}، ويدخل مربّعًا، أي \qty{6}{\percent}؛ والتصحيح \qty{3}{\percent}
يقع ضمن ذلك عدم اليقين.

\textbf{21.} ارتفاع حرارة حقل الفومارولات (الجزء~II) مع الزمن، وزيادة
في الزلازل الصغيرة تحت الانتفاخ، وتغيّر الغاز
المنبعث، ورفعٌ يُقاس بتحديد المواقع بالأقمار الصناعية أو بالتسوية عند
نقاط قليلة على سفح البركان.

\textbf{22.} نما الخزان تحت جبل هالفرين بمقدار
\textbf{15~مليون~\unit[text-series-to-math=true]{m^3}} ($\num{1.51e7}\,\unit{m^3}$).
\end{solution}
